% =========================================================
% 2. Create() ORM Method
% =========================================================
\section{\texttt{Create()} ORM Method}

\begin{frame}[standout]
  \texttt{Create()} ORM Method
\end{frame}

\begin{frame}[fragile]{\texttt{CREATE()} METHOD}
  \begin{itemize}
    \item The \texttt{create()} method is an ORM (Object Relational Mapping) method used to
          create one or more new records in an Odoo model.
    \item The \texttt{create()} method returns a \textbf{recordset} containing the newly created record.
  \end{itemize}
  \begin{block}{Method Syntax}
    The standard syntax of the \texttt{create()} method is:
\begin{lstlisting}
@api.model
def create(self, vals):
    return super().create(vals)
\end{lstlisting}
    This method has four important parts:
    \begin{enumerate}
      \item \texttt{@api.model}
      \item \texttt{create}
      \item \texttt{self}
      \item \texttt{vals}
    \end{enumerate}
  \end{block}
\end{frame}

\begin{frame}[fragile]{BREAKING DOWN THE METHOD SIGNATURE}
  \begin{itemize}
    \item \texttt{@api.model}: Indicates that the method is executed at the model level. The
          \texttt{@api.model} decorator tells Odoo that the method belongs to the model, not to
          an existing record.
    \item \texttt{create}: The ORM method responsible for creating records.
    \item \texttt{self}: The model (\texttt{student.student}), not a specific student record.
          If we are using \texttt{student.student} model, \texttt{self} does not represent a student record.
          Instead, \texttt{self} represents the model \texttt{student.student}.
  \end{itemize}
\begin{lstlisting}
print(self)
\end{lstlisting}
  Output is:
\begin{lstlisting}
student.student()
\end{lstlisting}
\end{frame}

\begin{frame}[fragile]{BREAKING DOWN THE METHOD SIGNATURE}
  \begin{itemize}
    \item \texttt{vals}: A dictionary containing the values to be stored in the database.
  \end{itemize}
\begin{lstlisting}
name = fields.Char()
birth_date = fields.Date()
gender = fields.Selection(...)
\end{lstlisting}
\begin{lstlisting}
vals = {
    'name': 'Ahmed Mohamed',
    'birth_date': date(2005, 5, 10),
    'gender': 'male',
}
\end{lstlisting}
  \begin{itemize}
    \item Each key corresponds to a field name in your model, and
    \item Each value is the data to store.
    \item The ORM maps each key to its corresponding field.
  \end{itemize}
\end{frame}

\begin{frame}[fragile]{WHAT DOES \texttt{CREATE()} RETURN?}
  \begin{itemize}
    \item The \texttt{create()} method returns a recordset.
  \end{itemize}
\begin{lstlisting}
student = self.env['student.student'].create({
    'name': 'Ahmed',
})
\end{lstlisting}
  \begin{itemize}
    \item Now, if we try to print the student using \texttt{print(student)}, we are getting this:
  \end{itemize}
\begin{lstlisting}
student.student(15,)
\end{lstlisting}
\end{frame}

\begin{frame}[fragile]{\texttt{@API.MODEL} VS \texttt{@API.MODEL\_CREATE\_MULTI}}
  There are two common decorators for \texttt{create()}.
  \begin{enumerate}
    \item \texttt{@api.model}: Used for creating a single record.
\begin{lstlisting}
@api.model
def create(self, vals):
\end{lstlisting}
          \texttt{vals} is a dictionary.
\begin{lstlisting}
student = self.env['student.student'].create({
    'name': 'Ahmed Mohamed',
})
\end{lstlisting}
    \item \texttt{@api.model\_create\_multi}: Used for creating multiple records efficiently.
\begin{lstlisting}
@api.model_create_multi
def create(self, vals_list):
\end{lstlisting}
  \end{enumerate}
\end{frame}

\begin{frame}[fragile]{\texttt{@API.MODEL} VS \texttt{@API.MODEL\_CREATE\_MULTI}}
  \begin{itemize}
    \item \texttt{vals\_list} is a list of dictionaries, i.e.,
  \end{itemize}
\begin{lstlisting}
[
    {'name': 'Ahmed'},
    {'name': 'Amina'},
]
\end{lstlisting}
  \textbf{Method Call}
  \begin{itemize}
    \item A method call is when we execute the method.
  \end{itemize}
\begin{lstlisting}
student = self.env['student.student'].create({
    'name': 'Ahmed Mohamed',
    'email': 'ahmed@example.com',
})
\end{lstlisting}
  \begin{itemize}
    \item Here, \texttt{.create(...)} is the method call.
  \end{itemize}
\end{frame}
